\chapter{Conclusion}%
\label{ch:conclusion}

\section{Summary \& Findings}

This thesis evaluated ResNet-18 inference as a staged microservice system. The study
started with controlled local split selection, moved the selected chain family into
local Kubernetes and AKS, and then measured the additional cost of service-mesh mTLS,
authorization policy, and AMD SEV-SNP confidential VM execution.

The main finding is that microservice inference cost is shaped by three linked factors:
where the neural network is split, where the resulting services are deployed, and how
strongly the communication and execution path are hardened. RQ1 showed that suitable
ResNet-18 boundaries lie in the mid-to-late region around \texttt{layer2} and
\texttt{layer3}, because early boundaries transfer larger activations and overly late
boundaries become compute-degenerate. RQ1.4 and RQ1.5 then showed that the same
service-count ordering remains visible in Kubernetes and AKS, although absolute latency
is environment-specific. RQ1.5b additionally bounded the inter-node penalty that the
single-node primary design intentionally hides: forcing every chain hop across a node
boundary on a multi-node AKS cluster adds between +2.131~ms and +4.086~ms per chained
condition, with a +2.305~ms baseline penalty on the monolithic configuration and the
architectural ordering preserved. RQ2 showed that mTLS/AuthZ and selective AMD SEV-SNP
confidential VM execution on the downstream service are feasible for the selected AKS
chain, but each protection layer adds measurable overhead; the security claims are
bounded by the mTLS termination boundary and by the SEV-SNP threat
model~\cite{he2021_attacking_collaborative,schluter2024_heckler}.

\section{Contributions}

The thesis makes three contributions. First, it provides a staged empirical evaluation
method for neural-network microservice decomposition, moving from local split selection
to Kubernetes, AKS, and security hardening without changing the model or benchmark
contract. Second, it shows that activation-transfer size and compute placement must be
considered together when selecting ResNet-18 service boundaries. Third, it quantifies a
security-performance trade-off for the selected Azure deployment path: mTLS/AuthZ is a
practical communication-hardening baseline, and selective AMD SEV-SNP TEE on the
downstream service adds an additional bounded cost while leaving the SEV-SNP residual
threat model unchanged. Extending the confidential-execution boundary to both services
in the chain is treated as future work rather than as a measured trade-off point.

\section{Limitations \& Threats to Validity}

The conclusions are bounded by the evaluated workload and deployment setup
(ResNet-18, CPU-only FP32 inference, batch size one, synchronous gRPC, and
same-node Kubernetes and AKS placement for the thesis-facing primary runs) and
by the bounded threat models of the evaluated security mechanisms. The
RQ1.5b sensitivity stage explicitly relaxes the same-node placement
constraint to bound the inter-node network penalty, but other production
factors (multi-region placement, autoscaling, concurrent client load,
non-default CNI plugins) remain out of scope. The full discussion of these
threats is consolidated in \cref{sec:analysis-threats}.

\section{Future Work}

Future work should extend the evaluation to larger and more diverse models, including
transformer-based inference workloads and models with different activation profiles.
The same staged method could also be applied under concurrent load, GPU execution,
cross-availability-zone or cross-region placement, and autoscaling conditions; the
multi-node sensitivity stage in \cref{sec:rq15b} bounds the intra-zone inter-node
penalty but leaves these wider topology questions open.
Another useful extension is to combine architectural partitioning with feature
compression, since compression may change the cost of early split points. Finally,
the security evaluation could be expanded to compare sidecar and sidecarless meshes,
application-layer encryption, attestation workflow integration, and different
confidential-computing platforms. In particular, extending the confidential-execution
boundary from the selective \texttt{service2\_only} scope evaluated here to both
services in the chain (full two-service TEE), and evaluating SEV-SNP residual
attack-surface attacks such as malicious-host
interrupts~\cite{schluter2024_heckler}, are explicit next steps that fall outside
the canonical evidence base of this thesis.
